\subsection{\texorpdfstring{THE ADDITIVE GROUP \(\mathbf{R}^n\)}{THE ADDITIVE GROUP R TO THE N}}

The set \(R^{n}\) , endowed with the group structure which is the product of the additive group structures of the n factors of \(R^{n}\) , is an abelian group; we use the additive notation, the sum of \(x = (x_{i})\) and \(y = (y_{i})\) being therefore \(x + y = (x_{i} + y_{i})\) . The topology of the number space is compatible with this group structure; endowed with these two structures, \(R^{n}\) is a topological group called the additive group of n-dimensional real number space. If n = 0, we make the convention that \(R^{0}\) denotes a group consisting only of the identity element.

The uniform structure of this group, called the additive uniformity of \(\mathbf{R}^n\) , is the product of the uniformities of the factors of \(\mathbf{R}^n\) (Chapter III, § 3, no. 2). If, for each integer \(p > 0\) , \(\mathrm{V}_p\) denotes the set of pairs \((x, y)\) of \(\mathbf{R}^n\) such that \(\max |x_i - y_i| \leqslant 1/p\) , the sets \(\mathrm{V}_p\) form a fundamental

\[
\mathbf {1} \leqslant i \leqslant n
\]

system of entourages of this uniformity. Whenever we consider \(R^{n}\) as a uniform space we shall always have in mind the additive uniformity just defined, unless the contrary is expressly stated. Endowed with this uniform structure, \(R^{n}\) is a complete uniform space (Chapter II, § 3, no. 5, Proposition 10).

